\documentclass[12pt]{article}
\usepackage[utf8]{inputenc}
\usepackage[T1]{fontenc}
\usepackage{hyperref}
\title{Lish\'an Noshan Language Resources for K-12 Curriculum Development}
\author{Compiled by Claude Sonnet 5 (claude-sonnet-5) for Andr\'as Kornai}
\date{}
\begin{document}
\maketitle
\section{Introduction}
Lish\'an Noshan is a Jewish Neo-Aramaic variety, now spoken mainly in Israel. Caution: my search found material on Lishanid Noshan (Inter-Zab Jewish Neo-Aramaic), whose ISO 639-3 code is given there as aij, not trg. I could not confirm that these are the same language as trg, so every item below is tied to that identification only provisionally. Wikipedia gives about 3,500 native speakers (2018), and the Endangered Languages Project lists the language as threatened.
\section{Primary Educational Resources}
I found no primers, textbooks or school programmes.
\section{Linguistic and Reference Materials}
\begin{itemize}
\item Overview of Inter-Zab Jewish Neo-Aramaic: \url{https://en.wikipedia.org/wiki/Inter-Zab_Jewish_Neo-Aramaic}
\item Geoffrey Khan, \emph{A Grammar of Neo-Aramaic: the dialect of the Jews of Arbel} (Brill, 1999). This is the main grammar of the Arbel dialect, mentioned in the Wikipedia article above. No separate URL was found.
\end{itemize}
Other works that appeared (Mutzafi on Koy Sanjaq, Sabar's dictionary of northwestern dialects) cover different dialects, so I do not list them as resources for this language.
\section{Religious and Community Resources}
I found none.
\section{Digital and Audio Resources}
\begin{itemize}
\item Endangered Languages Project, resources page, with video and documentation items: \url{https://endangeredlanguages.com/node/39005/resources}
\item Endangered Languages Project, language entry: \url{https://endangeredlanguages.com/lang/3697?hl=zh-cn}
\end{itemize}
\section{Curriculum Development Context}
The language is not known to be taught in schools, and its speakers are mostly older adults. The identification question above should be settled before any curriculum work begins. Documentation of the Arbel dialect, such as Khan's grammar, is the most substantial scholarly basis found.
\end{document}